\begin{thebibliography}{99}
\small

\bibitem{vina2010} Trott O, Olson AJ (2010) AutoDock Vina: improving the speed and
accuracy of docking with a new scoring function, efficient optimization, and
multithreading. J Comput Chem 31:455--461. doi:10.1002/jcc.21334.

\bibitem{ding2023} Ding J, Tang S, Mei Z et al (2023) Vina-GPU 2.0: further
accelerating AutoDock Vina and its derivatives with graphics processing units.
J Chem Inf Model 63:1982--1998. doi:10.1021/acs.jcim.2c01504.

\bibitem{dockstring2022} Garc\'ia-Orteg\'on M, Simm GNC, Tripp AJ,
Hern\'andez-Lobato JM, Bender A, Bacallado S (2022) DOCKSTRING: easy molecular
docking yields better benchmarks for ligand design. J Chem Inf Model
62:3486--3502. doi:10.1021/acs.jcim.1c01334.

\bibitem{nikitin2025} Nikitin I, Morgunov I, Safronov V, Kalyuzhnaya A,
Fedorov M (2025) Towards explainable computational toxicology: linking
antitargets to rodent acute toxicity. Pharmaceutics 17:1573.
doi:10.3390/pharmaceutics17121573.

\bibitem{pan2003} Pan Y, Huang N, Cho S, MacKerell AD Jr (2003) Consideration of
molecular weight during compound selection in virtual target-based database
screening. J Chem Inf Comput Sci 43:267--272. doi:10.1021/ci020055f.

\bibitem{vigers2004} Vigers GPA, Rizzi JP (2004) Multiple active site corrections
for docking and virtual screening. J Med Chem 47:80--89. doi:10.1021/jm030161o.

\bibitem{jacobsson2006} Jacobsson M, Karl\'en A (2006) Ligand bias of scoring
functions in structure-based virtual screening. J Chem Inf Model 46:1334--1343.
doi:10.1021/ci050407t.

\bibitem{carta2007} Carta G, Knox AJS, Lloyd DG (2007) Unbiasing scoring
functions: a new normalization and rescoring strategy. J Chem Inf Model
47:1564--1571. doi:10.1021/ci600471m.

\bibitem{sepresa2009} Yang L, Luo H, Chen J, Xing Q, He L (2009) SePreSA: a
server for the prediction of populations susceptible to serious adverse drug
reactions implementing the methodology of a chemical--protein interactome.
Nucleic Acids Res 37:W406--W412. doi:10.1093/nar/gkp312.

\bibitem{casey2009} Casey FP, Pihan E, Shields DC (2009) Discovery of small
molecule inhibitors of protein--protein interactions using combined ligand and
target score normalization. J Chem Inf Model 49:2708--2717.
doi:10.1021/ci900294x.

\bibitem{luo2017} Luo Q, Zhao L, Hu J, Jin H, Liu Z, Zhang L (2017) The scoring
bias in reverse docking and the score normalization strategy to improve success
rate of target fishing. PLoS One 12:e0171433.
doi:10.1371/journal.pone.0171433.

\bibitem{lapillo2019} Lapillo M, Tuccinardi T, Martinelli A, Macchia M,
Giordano A, Poli G (2019) Extensive reliability evaluation of docking-based
target-fishing strategies. Int J Mol Sci 20:1023.
doi:10.3390/ijms20051023.

\bibitem{lee2020crds} Lee A, Kim D (2020) CRDS: consensus reverse docking system
for target fishing. Bioinformatics 36:959--960.
doi:10.1093/bioinformatics/btz656.

\bibitem{krause2023} Krause F, Voigt K, Di Ventura B, {\"O}zt{\"u}rk MA (2023)
ReverseDock: a web server for blind docking of a single ligand to multiple
protein targets using AutoDock Vina. Front Mol Biosci 10:1243970.
doi:10.3389/fmolb.2023.1243970.

\bibitem{fukunishi2005} Fukunishi Y, Mikami Y, Nakamura H (2005) Similarities
among receptor pockets and among compounds: analysis and application to
in silico ligand screening. J Mol Graph Model 24:34--45.
doi:10.1016/j.jmgm.2005.04.004.

\bibitem{fukunishi2006} Fukunishi Y, Mikami Y, Takedomi K, Yamanouchi M,
Shima H, Nakamura H (2006) Classification of chemical compounds by
protein--compound docking for use in designing a focused library. J Med Chem
49:523--533. doi:10.1021/jm050480a.

\bibitem{fukunishi2018} Fukunishi Y, Yamashita Y, Mashimo T, Nakamura H (2018)
Prediction of protein--compound binding energies from known activity data:
docking-score-based method and its applications. Mol Inform 37:e1700120.
doi:10.1002/minf.201700120.

\bibitem{lee2012profiles} Lee M, Kim D (2012) Large-scale reverse docking
profiles and their applications. BMC Bioinformatics 13(Suppl 17):S6.
doi:10.1186/1471-2105-13-S17-S6.

\bibitem{martin2011pqsar} Martin E, Mukherjee P, Sullivan D, Jansen J (2011)
Profile-QSAR: a novel meta-QSAR method that combines activities across the
kinase family to accurately predict affinity, selectivity, and cellular
activity. J Chem Inf Model 51:1942--1956. doi:10.1021/ci1005004.

\bibitem{bembenek2018} Bembenek SD, Hirst G, Mirzadegan T (2018)
Determination of a focused mini kinase panel for early identification of
selective kinase inhibitors. J Chem Inf Model 58:1434--1440.
doi:10.1021/acs.jcim.8b00222.

\bibitem{zhang2019informer} Zhang H, Ericksen SS, Lee CP et al (2019)
Predicting kinase inhibitors using bioactivity matrix derived informer sets.
PLoS Comput Biol 15:e1006813. doi:10.1371/journal.pcbi.1006813.

\bibitem{laufkotter2020} Laufk\"otter O, Laufer S, Bajorath J (2020)
Identifying representative kinases for inhibitor evaluation via systematic
analysis of compound-based target relationships. Eur J Med Chem 204:112641.
doi:10.1016/j.ejmech.2020.112641.

\bibitem{mangione2019cando} Mangione W, Samudrala R (2019) Identifying protein
features responsible for improved drug repurposing accuracies using the CANDO
platform: implications for drug design. Molecules 24:167.
doi:10.3390/molecules24010167.

\bibitem{kangas2014active} Kangas JD, Naik AW, Murphy RF (2014) Efficient
discovery of responses of proteins to compounds using active learning. BMC
Bioinformatics 15:143. doi:10.1186/1471-2105-15-143.

\bibitem{davis2011} Davis MI, Hunt JP, Herrgard S et al (2011) Comprehensive
analysis of kinase inhibitor selectivity. Nat Biotechnol 29:1046--1051.
doi:10.1038/nbt.1990.

\bibitem{wu2025davis} Wu MH (2025) A Complete and Modification-Aware DAVIS
Dataset. Harvard Dataverse, V3. doi:10.7910/DVN/RTQGP1.

\bibitem{drewry2017} Drewry DH, Wells CI, Andrews DM et al (2017) Progress
towards a public chemogenomic set for protein kinases and a call for
contributions. PLoS One 12:e0181585. doi:10.1371/journal.pone.0181585.

\bibitem{anastassiadis2011} Anastassiadis T, Deacon SW, Devarajan K,
Ma H, Peterson JR (2011) Comprehensive assay of kinase catalytic activity
reveals features of kinase inhibitor selectivity. Nat Biotechnol 29:1039--1045.
doi:10.1038/nbt.2017.

\bibitem{mazzucato2016} Mazzucato L, Fontanini A, La Camera G (2016) Stimuli
reduce the dimensionality of cortical activity. Front Syst Neurosci 10:11.
doi:10.3389/fnsys.2016.00011.

\bibitem{roy2007} Roy O, Vetterli M (2007) The effective rank: a measure of
effective dimensionality. In: 15th European Signal Processing Conference,
Pozna\'n, pp 606--610. doi:10.5281/zenodo.40328.

\bibitem{hubert1976} Hubert L, Schultz J (1976) Quadratic assignment as a
general data analysis strategy. Br J Math Stat Psychol 29:190--241.
doi:10.1111/j.2044-8317.1976.tb00714.x.

\bibitem{chen2010} Chen Y, Wiesel A, Eldar YC, Hero AO (2010) Shrinkage
algorithms for MMSE covariance estimation. IEEE Trans Signal Process
58:5016--5029. doi:10.1109/TSP.2010.2053029.

\bibitem{murcko1996} Bemis GW, Murcko MA (1996) The properties of known drugs.
1. Molecular frameworks. J Med Chem 39:2887--2893.
doi:10.1021/jm9602928.

\bibitem{morgan1965} Morgan HL (1965) The generation of a unique machine
description for chemical structures---a technique developed at Chemical
Abstracts Service. J Chem Doc 5:107--113. doi:10.1021/c160017a018.

\bibitem{rogers2010} Rogers D, Hahn M (2010) Extended-connectivity fingerprints.
J Chem Inf Model 50:742--754. doi:10.1021/ci100050t.

\bibitem{butina1999} Butina D (1999) Unsupervised data base clustering based on
Daylight's fingerprint and Tanimoto similarity: a fast and automated way to
cluster small and large data sets. J Chem Inf Comput Sci 39:747--750.
doi:10.1021/ci9803381.

\bibitem{rdkit2025} Landrum G et al (2025) RDKit: open-source cheminformatics
software, release 2025.03.6. Zenodo. doi:10.5281/zenodo.16996017.

\bibitem{numpy2020} Harris CR, Millman KJ, van der Walt SJ et al (2020) Array
programming with NumPy. Nature 585:357--362.
doi:10.1038/s41586-020-2649-2.

\bibitem{scipy2020} Virtanen P, Gommers R, Oliphant TE et al (2020) SciPy 1.0:
fundamental algorithms for scientific computing in Python. Nat Methods
17:261--272. doi:10.1038/s41592-019-0686-2.

\bibitem{matplotlib2007} Hunter JD (2007) Matplotlib: a 2D graphics environment.
Comput Sci Eng 9:90--95. doi:10.1109/MCSE.2007.55.

\end{thebibliography}
